% !TeX root = main.tex
%!TEX root = main.tex





\section{The Blueprint Construction}\label{sec:construct-blueprint}

We provide our blueprint construction in this section, proving the following theorem. 
\begin{theorem*}[Restatement of~\Cref{thm:blueprint}]
	For any sufficiently large integer $m$, there exists a proper blueprint $\bG_m$ with parameters $P=\poly(m)$ and $C=m^{O(m)}$ such that 
	\[
		value(\bG_m) = \frac{2}{3} - \Theta\!\paren{\frac{1}{m}}. 
	\]
\end{theorem*}

We start by presenting the setup needed for this construction, followed by the construction itself and its properness proof, and finally calculate its value to conclude the proof of~\Cref{thm:blueprint}. 

\subsection{Setup} 
Let $m$ be a large integer that is used to define all our other parameters for the blueprint. We define a blueprint $\bG := \bG(m)$ in this section. 
Define: 
\begin{align*}
L &:= m^4, \tag{used to determine the \emph{length} of random walks} \\
D &:= 2m \cdot L, \tag{used to determine the (range of) \emph{delays} in random walks} \\
P &:= L+D+1 \tag{parameter $P$ of the blueprint} \\
C &:= \text{sufficiently large as function of $m$ to be fixed later} \tag{parameter $C$ of the blueprint}. 
\end{align*}
Recall that vertices of $\bG$ are $\bL,\bR$ and they are labeled by $[C]^{P}$. We partition indices of $P$ as follows: 
\begin{gather*}
	P := (\pstart , \pdelay , \pwalk) \quad \text{where} \\ 
	\pstart=1, \quad \pdelay=2, \quad \text{and,} \quad {\pwalk} = [3: P]. 
\end{gather*}

We use these indices to associate a delayed random walk to vertices of the blueprint as follows. 

\begin{definition}\label{def:W(x)}
Sample $x \in [C]^P$ uniformly at random and define a \textbf{delayed random walk} $W = W(x) = w_1,w_2,\cdots,w_T$ for some $T \leq D+L$ as follows\footnote{It is possible that the walk does not terminate and in that case we only consider its $D+L$ first states.}.
\begin{enumerate}[itemsep=3pt]
\item\label{line:pstart} $x_{\pstart}$ determines the start state $w_1$ of the random walk $W$: we partition $[C]$ into intervals of proportional length $\mu(1),\mu(2),\ldots,\mu(m)$ (as defined in~\Cref{lem:start-state}), 
and pick $w_1$ depending on which interval $x_{\pstart}$ lands in. 
\item\label{line:pdelay} $x_{\pdelay}$ determines the delay of the delayed random walk $W$: we partition $[C]$ into $D$ equal-length intervals and the delay of the random walk $d(x)=d$
if $x_{\pdelay}$ lands in $d$-th interval.  
\item\label{line:reflection} $x_{\pwalk}$ determines the inner walk: for any time $t \geq d(x)$, $x_{t+2}$ determines the transition from $w_t$ to $w_{t+1}$ in $W$, where
$w_{t+1} = w_t -1$ if $x_{t+2} \in [1:{C}/{2}]$ and $w_{t+1}=w_t + 1$ otherwise.
The walk terminates whenever it hits $0$ or $2m$ and the subsequent states remain the same. 
\end{enumerate}

We also define the \textbf{reflection walk} of $x$, denoted by $\barW(x)$, exactly as $W(x)$ except for Line~\eqref{line:reflection}: we now increase the state if $x_i \in [1:C/2]$ and decrease it otherwise (the exact opposite rule as in the definition of $W$).  
\end{definition}
See~\Cref{fig:W(x)} for an illustration of this definition and further help with parsing it. 

\clearpage
Before getting to the main part, a quick remark is in order. 
\begin{remark}
We pick $C$ such that 
\begin{itemize}
	\item $\mu(i) \cdot C$ is an integer for all $i \in [2m-1]$ (so the distribution of the start point in Line~\eqref{line:pstart} 
	matches $\mu$);  
	\item $D | C$ (so the distribution of the delay in Line~\eqref{line:pdelay} is uniform over $[D]$); 
	\item $C$ is even (so the distribution of the inner walk in Line~\eqref{line:reflection} is a standard random walk).  
\end{itemize}
	We can thus set 
	\[
		C = 2D \cdot \prod_{i=1}^{m} \textnormal{denominator of $\mu(i)$}
	\]
	and given the values of parameters $\mu(i)$ in~\Cref{lem:start-state} and since $D= \poly(m)$, we have $C=m^{O(m)}$. 
\end{remark}


\begin{figure}[t!]
\centering
\begin{tikzpicture}[
    x=0.78cm,
    y=1cm,
    >=Latex,
    line cap=butt
]

    %------------------------------------------------------------
    % Geometry
    %------------------------------------------------------------
    \def\boxheight{2.8}
    \def\boxwidth{16}

    % End of the first and second boxes
    \def\startend{1.6}
    \def\delayend{3.6}

    % Partitions inside the walk region
    \def\walkA{6.4}
    \def\walkB{7.55}
    \def\walkC{8.70}
    \def\walkD{9.85}

    % Half-length of each small horizontal interval separator
    \def\separatorhalfwidth{0.12}

    % Vertical midpoint of the large box
    \pgfmathsetmacro{\ymid}{\boxheight/2}

    % Centers of the first two boxes
    \pgfmathsetmacro{\startcenter}{\startend/2}
    \pgfmathsetmacro{\delaycenter}{(\startend+\delayend)/2}

    % Centers of the three narrow boxes
    \pgfmathsetmacro{\dcenterA}{(\walkA+\walkB)/2}
    \pgfmathsetmacro{\dcenterB}{(\walkB+\walkC)/2}
    \pgfmathsetmacro{\dcenterC}{(\walkC+\walkD)/2}

    % The walk region is everything after the delay box
    \pgfmathsetmacro{\walkcenter}{(\delayend+\boxwidth)/2}

    %------------------------------------------------------------
    % Styles
    %------------------------------------------------------------
    \tikzset{
        outer box/.style={
            draw=black,
            rounded corners=2mm,
            line width=0.65pt
        },
        partition/.style={
            draw=black,
            line width=0.55pt
        },
        colored bar/.style={
            line width=3pt,
            line cap=butt
        },
        interval separator/.style={
            draw=black,
            line width=0.55pt,
            line cap=round
        },
        small arrow/.style={
            -{Latex[length=1.8mm,width=1.35mm]},
            line width=0.9pt,
            black
        },
        upward arrow/.style={
            -{Latex[length=2.7mm,width=2.1mm]},
            line width=1.5pt,
            green!65!black,
            line cap=butt
        },
        downward arrow/.style={
            -{Latex[length=2.7mm,width=2.1mm]},
            line width=1.5pt,
            red!70,
            line cap=butt
        },
        lower label/.style={
            anchor=north,
            font=\large\itshape
        },
        set label/.style={
            font=\Large
        }
    }

    %------------------------------------------------------------
    % Clip all interior elements to the rounded outer box
    %------------------------------------------------------------
    \begin{scope}
        \clip[rounded corners=2mm]
            (0,0) rectangle (\boxwidth,\boxheight);

        %--------------------------------------------------------
        % Exact top-to-bottom partition lines
        %--------------------------------------------------------
        \foreach \x in {
            \startend,
            \delayend,
            \walkA,
            \walkB,
            \walkC,
            \walkD
        }{
            \draw[partition]
                (\x,0)
                --
                (\x,\boxheight);
        }

        %--------------------------------------------------------
        % First box: unequal interval lengths
        %--------------------------------------------------------
        \draw[colored bar,green!65!black]
            (\startcenter,0)
            --
            (\startcenter,0.55);

        \draw[colored bar,red!75]
            (\startcenter,0.55)
            --
            (\startcenter,1.90);

        \draw[colored bar,blue!60]
            (\startcenter,1.90)
            --
            (\startcenter,\boxheight);

        % Horizontal separators between the three intervals
        \foreach \y in {0.55,1.90}{
            \draw[interval separator]
                ({\startcenter-\separatorhalfwidth},\y)
                --
                ({\startcenter+\separatorhalfwidth},\y);
        }

        %--------------------------------------------------------
        % Second box: six equal intervals
        %--------------------------------------------------------
        \foreach \col [count=\i from 0] in {
            red!75,
            orange!85!black,
            cyan!65!black,
            green!60!black,
            magenta!65,
            blue!60
        }{
            \pgfmathsetmacro{\segmentbottom}{\i*\boxheight/6}
            \pgfmathsetmacro{\segmenttop}{(\i+1)*\boxheight/6}

            \draw[colored bar,\col]
                (\delaycenter,\segmentbottom)
                --
                (\delaycenter,\segmenttop);
        }

        % Horizontal separators between the six equal intervals
        \foreach \i in {1,...,5}{
            \pgfmathsetmacro{\separatorheight}
                {\i*\boxheight/6}

            \draw[interval separator]
                ({\delaycenter-\separatorhalfwidth},
                 \separatorheight)
                --
                ({\delaycenter+\separatorhalfwidth},
                 \separatorheight);
        }

        %--------------------------------------------------------
        % Three narrow boxes:
        % each is divided into equal upper and lower intervals
        %--------------------------------------------------------
        \foreach \x in {
            \dcenterA,
            \dcenterB,
            \dcenterC
        }{
            % Upper interval
            \draw[upward arrow]
                (\x,\ymid)
                --
                (\x,\boxheight);

            % Lower interval
            \draw[downward arrow]
                (\x,\ymid)
                --
                (\x,0);

            % Horizontal separator between the two intervals
            \draw[interval separator]
                ({\x-\separatorhalfwidth},\ymid)
                --
                ({\x+\separatorhalfwidth},\ymid);
        }

        %--------------------------------------------------------
        % Small left-to-right arrows
        %--------------------------------------------------------

        % First box
        \draw[small arrow]
            (0.32,1.23)
            --
            (0.59,1.23);

        % Second box
        \draw[small arrow]
            (2.08,2.15)
            --
            (2.35,2.15);

        % First narrow box
        \draw[small arrow]
            (6.57,2.02)
            --
            (6.82,2.02);

        % Second narrow box
        \draw[small arrow]
            (7.70,0.70)
            --
            (7.95,0.70);

        % Third narrow box
        \draw[small arrow]
            (8.85,0.95)
            --
            (9.10,0.95);
    \end{scope}

    %------------------------------------------------------------
    % Draw the outer boundary last
    %------------------------------------------------------------
    \draw[outer box]
        (0,0)
        rectangle
        (\boxwidth,\boxheight);

    %------------------------------------------------------------
    % External labels
    %------------------------------------------------------------
\node[set label,overlay]
    at (-0.72,\ymid)
    {$[C]$};

    \node[set label]
        at ({\boxwidth/2},{\boxheight+0.62})
        {$[P]$};

    % All lower labels share the same baseline
    \node[lower label]
        at (\startcenter,-0.1)
        {\pstart};

    \node[lower label]
        at (\delaycenter,-0.1)
        {\pdelay};

    \node[lower label]
        at (\dcenterA,-0.1)
        {$d+2$};

    \node[lower label]
        at (\walkcenter,-0.1)
        {\pwalk};

\end{tikzpicture}

\caption{An illustration of the random walk $W(x)$ in~\Cref{def:W(x)} when $m=3$ and $D=6$. The figure shows how the indices of $[P]$ are partitioned into $(\pstart,\pdelay,\pwalk)$. The intervals inside each box shows how $[C]$ is partitioned into different intervals. The tiny arrows in each box shows the value of a single point $x \in [C]^P$ on each relevant index. The walk $W(x)$ thus first considers $x_{\pstart}$ and so chooses starting state $2$. The delay $d$ is set to $5$ given $x_{\pdelay}$. Finally, the walk waits until time $d$ and then decide its next move based on the index $d+2$ and so on; in this particular example, the walk goes to $3$ next, then $2$, and then $1$. }
\label{fig:W(x)}
\end{figure}


We use walks $W(x)$ for defining subsets of $\bL$ and reflection walks $\barW(x)$ for subsets in $\bR$. Specifically, for any integer $j \in [2m-1]$ and time $t \in [D+L]$, we define
\begin{align*}
	\bL_{j,t} &:= \set{\bL(x) : \text{$W(x)$ is $L$-restricted, $W(x)$ reaches $0$, $\maxwalk{W(x)}=j$, and $\timemax{W(x)} = t$}}, \\
	\bR_{j,t} &:= \set{\bR(x) : \text{$\barW(x)$ is $L$-restricted, $\barW(x)$ reaches $0$, $\maxwalk{\barW(x)}=j$, and $\timemax{\barW(x)} = t$}}. 
\end{align*}
These sets will play the key role in the definition of edges of the blueprint (they are essentially the endpoints of the edges). 

We have the following immediate relation between the sizes of these sets. 

\begin{claim}\label{obs:bjt-direct}
	For all $j \in [2m-1]$ and $t \in [D+L]$, $\card{\bL_{j,t}} = \card{\bR_{j,t}}$.
\end{claim}
\begin{proof}
	The only difference between $\bL_{j,t}$ and $\bR_{j,t}$ is in Line~\eqref{line:reflection} of~\Cref{def:W(x)}: we thus have $\bL(x) \in \bL_{j,t}$ iff $\bR(\overline{x}) \in \bR_{j,t}$, where
	$\overline{x}$ is defined to be the same as $x$ in $[P] \setminus \pwalk$ and for $i \in \pwalk$, $\overline{x}_i := C - x_i + 1$. Using the bijection between $x$'s and $\overline{x}$'s, we can conclude the proof. 
\end{proof}

The more interesting property is the near-equality of the sizes of $\bL_{j,t}$ and $\bR_{2m-j,t}$ (whenever $t \in [L:D]$); this is where we use most of the properties we incorporated in our random walks. 
\begin{lemma}\label{lem:bjt-reflect}
	For all $j \in [2m-1]$ and $t \in [L:D]$, 
	\[
	\card{\bR_{2m-j,t}} - 2^{-m} \cdot \frac{C^P}{D} \leq \card{\bL_{j,t}} \leq \card{\bR_{2m-j,t}} + 2^{-m} \cdot \frac{C^P}{D}. 
	\]
\end{lemma}
\begin{proof}
	Define $W_{in}(x)$ as the inner walk of $W(x)$, which is the walk that starts at time $d(x)$ (the delay) and continues the same as $W$; we define $\barW_{in}(x)$ similarly. 
	In the following, when clear from the context, we drop $(x)$ and write $W$ to denote $W(x)$ (similarly for $W_{in},\barW$, etc.). 
	
	We first have 
	\begin{align*}
		\frac{\card{\bL_{j,t}}}{C^P} &= \Pr_{x}\Paren{\text{$W$ $L$-restricted, $W$ reaches $0$, $\maxwalk{W}=j$, $\timemax{W} = t$}} \tag{by~\Cref{def:W(x)}}\\
		&=  \Pr_{x}\Paren{\text{$W_{in}$ $L$-restricted, $W_{in}$ reaches $0$, $\maxwalk{W_{in}}=j$,  $\timemax{W} = t$}} \tag{$L$-restrictedness, reaching $0$, and the max-states are all the same 
		between $W_{in}$ and $W$}  \\
		&= \Pr_{x}\Paren{\text{$W_{in}$ $L$-restricted, $W_{in}$ reaches $0$, $\maxwalk{W_{in}}=j$, $\timemax{W_{in}} = t-(d(x)-1)$}} \tag{by the definition of $\timemax{W}$ in~\Cref{eq:timemaxW}} \\
		&= \frac1D \cdot \Pr_{x}\Paren{\text{$W_{in}$ $L$-restricted, $W_{in}$ reaches $0$, $\maxwalk{W_{in}}=j$}} 
	\end{align*}
	as $d(x)$ is independent of $W_{in}$ and is uniform over $[D]$, 
	and so the event $d(x) = t-\timemax{W_{in}}+1$ (conditioned on $W_{in}$) happens with probability $1/D$ as the RHS is in $[D]$ as well; this is because $\timemax{W_{in}} \leq L$ (by $L$-restrictedness) and $t \in [L:D]$. 
	
	We can do the same exact analysis for $\bR_{2m-j,t}$ and $\barW$ as well to have 
	\[
		\frac{\card{\bR_{2m-j,t}}}{C^P} = \frac{1}{D} \cdot \Pr_{x}\Paren{\text{$\barW_{in}$ $L$-restricted, $\barW_{in}$ reaches $0$, $\maxwalk{\barW_{in}}=2m-j$}}. 
	\]
	Consider the expression in the RHS of either of these two terms. The distribution of $W_{in}$ or $\barW_{in}$ is exactly the same as the distribution of a random walk $\widetilde{W} \sim \dist_m$ defined in~\Cref{lem:start-state}. Hence,  
	\begin{align*}
		\frac{\card{\bL_{j,t}}}{C^P} &=  \frac{1}{D} \cdot \Pr_{\widetilde{W}}\Paren{\text{$\widetilde{W}$ $L$-restricted, $\widetilde{W}$ reaches $0$, $\maxwalk{\widetilde{W}}=j$}} \\
		&\geq \frac{1}{D} \cdot \paren{\Pr_{\widetilde{W}}\Paren{\text{$\widetilde{W}$ reaches $0$, $\maxwalk{\widetilde{W}}=j$}} - \Pr_{\widetilde{W}}\Paren{\text{$\widetilde{W}$ not $L$-restricted}}} \tag{as $\Pr(A \wedge B) \geq \Pr(B) - \Pr(\overline{A})$} \\
		&\geq \frac{1}{D} \cdot \paren{\Pr_{\widetilde{W}}\Paren{\text{$\widetilde{W}$ reaches $0$, $\maxwalk{\widetilde{W}}=j$}} - 2^{-m}} \tag{as $L=m^4$ and by~\Cref{clm:walk-length}} \\ 
		&= \frac{1}{D} \cdot \paren{\Pr_{\widetilde{W}}\Paren{\text{$\widetilde{W}$ reaches $0$, $\maxwalk{\widetilde{W}}=2m-j$}} - 2^{-m}} \tag{by the symmetry property of~\Cref{lem:start-state}} \\
		&\geq \frac{1}{D} \cdot \paren{\Pr_{\widetilde{W}}\Paren{\text{$\widetilde{W}$ $L$-restricted, $\widetilde{W}$ reaches $0$, $\maxwalk{\widetilde{W}}=2m-j$}} - 2^{-m}} \tag{as $\Pr(A) \geq \Pr(A \wedge B)$} \\
		&= \frac{\card{\bR_{2m-j,t}}}{C^P} - \frac{2^{-m}}{D}. 
	\end{align*}
	Applying the same argument starting from $\bR_{2m-j,t}$ then implies that 
	\[
		\card{\bR_{2m-j,t}} - 2^{-m} \cdot \frac{C^P}{D} \leq \card{\bL_{j,t}} \leq \card{\bR_{2m-j,t}} + 2^{-m} \cdot \frac{C^P}{D} \qed
	\]
	
\end{proof}

\subsection{The Construction} 

We are now ready to define the edges of the blueprint $\bG=(\bL,\bR,\bE)$ and prove its properness. As stated earlier, both vertex-sets $\bL$ and $\bR$ are labeled by $[C]^P$. The edges in $\bE$ are defined as follows: 
\begin{itemize}
	\item For any $j \in [2m-1]$ and $t \in [L:D]$, we add a matching $\bM_{j,t}$ of size $s_{j,t} := \min\paren{\card{\bL_{j,t}},\card{\bR_{2m-j,t}}}$ between two arbitrarily-chosen $s_{j,t}$-subsets of $\bL_{j,t}$ and $\bR_{2m-j,t}$. Moreover, 
	the matching $\bM_{j,t}$ is added as part of $\bEE{p}$ for $p=t+2$.
\end{itemize}

We need to check that the blueprint $\bG$ is proper, which is handled by the following two claims. 

\begin{claim}[``Matching constraint'']\label{clm:bG-matching}
	The edge-set $\bE$ is a matching. 
\end{claim}
\begin{proof}
	By definition, the sets $\bL_{j,t}$ are disjoint for different values of $j$ (the maximum of a random walk is a unique state) as well as $t$ (the \emph{first} time reaching the eventual maximum is unique). Same is true for $\bR_{2m-j,t}$ as well, 
	thus $\bE = \cup_{p \in [P]} \bEE{p}$ itself is a matching also.  
\end{proof}

% In the preamble:
% \usepackage{tikz}
% \usetikzlibrary{arrows.meta}

\begin{figure}[t!]
\centering
\begin{tikzpicture}[x=1.15cm,y=1cm,>=Latex]

    %------------------------------------------------------------
    % Parameters
    %------------------------------------------------------------
    \def\axisy{0}
    \def\topy{3.15}
    \def\texty{3.55}
    \def\circley{0.72}
    \def\dy{0.23}

    % Arrow levels, with fixed vertical spacing
    \pgfmathsetmacro{\ya}{2.30}
    \pgfmathsetmacro{\yb}{\ya-\dy}
    \pgfmathsetmacro{\yc}{\yb-\dy}
    \pgfmathsetmacro{\yd}{\yc-\dy}
    \pgfmathsetmacro{\ye}{\yd-\dy}
    \pgfmathsetmacro{\yf}{\ye-\dy}

    %------------------------------------------------------------
    % Styles
    %------------------------------------------------------------
    \tikzset{
        gridline/.style={
            dashed,
            draw=black!55,
            line width=0.5pt
        },
        baseline/.style={
            draw=black,
            line width=2.2pt
        },
        redarrow/.style={
            -{Latex[length=2.8mm,width=2.2mm]},
            draw=red!75!black,
            line width=1.3pt,
            shorten <=0pt,
            shorten >=0pt
        },
        greenarrow/.style={
            -{Latex[length=2.8mm,width=2.2mm]},
            draw=green!60!black,
            line width=1.3pt,
            shorten <=0pt,
            shorten >=0pt
        },
        ball/.style={
            circle,
            draw=black,
            fill=yellow!75,
            minimum size=5mm,
            inner sep=0pt
        }
    }

    %------------------------------------------------------------
    % Dashed vertical lines at every integer
    %------------------------------------------------------------
    \foreach \i in {0,...,12}{
        \draw[gridline]
            (\i,\axisy)
            --
            (\i,\topy);
    }

    %------------------------------------------------------------
    % Horizontal baseline
    %------------------------------------------------------------
    \draw[baseline]
        (0,\axisy)
        --
        (12,\axisy);

    %------------------------------------------------------------
    % Integer labels
    %------------------------------------------------------------
    \foreach \i in {0,...,12}{
        \node[font=\Large]
            at (\i,-0.55)
            {\i};
    }

    %------------------------------------------------------------
    % Circles and upper labels
    %------------------------------------------------------------
    \node[ball] at (2,\circley) {};
    \node[ball] at (10,\circley) {};

    \node[font=\Large]
        at (2,\texty)
        {$W(x_{<p}\circ z)$};

    \node[font=\Large]
        at (10,\texty)
        {$\barW(y_{<p}\circ z)$};

    %------------------------------------------------------------
    % Left side
    % Every arrow has length exactly one grid interval.
    %------------------------------------------------------------

    % Red arrows point left
    \draw[redarrow] (1,\ya) -- (0,\ya);
    \draw[redarrow] (2,\yb) -- (1,\yb);
    \draw[redarrow] (3,\yc) -- (2,\yc);
    \draw[redarrow] (4,\yd) -- (3,\yd);

    % Green arrows point right
    \draw[greenarrow] (2,\ye) -- (3,\ye);
    \draw[greenarrow] (3,\yf) -- (4,\yf);

    %------------------------------------------------------------
    % Right side
    % Reflected through x=6, with colors exchanged.
    %------------------------------------------------------------

    % Green arrows point right
    \draw[greenarrow] (8,\yd) -- (9,\yd);
    \draw[greenarrow] (9,\yc) -- (10,\yc);
    \draw[greenarrow] (10,\yb) -- (11,\yb);
    \draw[greenarrow] (11,\ya) -- (12,\ya);

    % Red arrows point left
    \draw[redarrow] (10,\ye) -- (9,\ye);
    \draw[redarrow] (9,\yf) -- (8,\yf);

\end{tikzpicture}

\caption{An illustration of the proof of the ban constraint (\Cref{clm:bG-ban}) for $\bG$ with $m=6$, $j=2$, and $z \approx (\uparrow,\uparrow,\downarrow,\downarrow,\downarrow,\downarrow)$, which by that we mean the $\uparrow$ indices are in $[(C/2)+1:C]$
and $\downarrow$ ones are in $[1:C/2]$. Suppose the choices of $x_{<p}$ and $y_{<p}$ already ``bring'' the walks $W(x_{<p} \circ z)$ and $\barW(y_{<p} \circ z)$ to 
states $2$ and $10$, respectively. From there, $W(x_{<p} \circ z)$, increases the state by two and then decreases it by four, whereas by reflection, $\barW(y_{<p} \circ z)$ does the exact opposite; thus, 
$W(x_{<p} \circ z)$ reaches state $0$ iff $\barW(y_{<p} \circ z)$ reaches $12 (=2m)$. As such, at most one pair of these vertices can be matched in $\bG$ (only vertices whose walk reaches $0$ are matched regardless of being in $\bL$ or $\bR$).}
\label{fig:W}
\end{figure}

\begin{claim}[``Ban constraint'']\label{clm:bG-ban}
	The edge-set $\bE$ respects the ban constraint. 
\end{claim}
\begin{proof}

	Consider an edge $\be \in \bEE{p}$ for  $p=t+2$ for some arbitrary $t \in [L:D]$ and suppose $\be$ is between some $\bL(x) \in \bL_{j,t}$ and $\bR(y) \in \bR_{2m-j,t}$ for $x,y \in [C]^P$. Recall all edges of $\bG$ are of this form. 
	
	By~\Cref{def:proper-blueprint}, the ban constraint requires that for all $z \in [C]^{P-p+1}$, at most one of $\bL(x_{<p} \circ z)$ and $\bR(y_{<p} \circ z)$ is matched in $\bE$. We show that this is the case for every fixed $z \in [C]^{P-p+1}$. 
	\Cref{fig:W} can be helpful with providing an illustration of this step of the proof. 
	


	
	Consider the walk $W(x_{<p} \circ z)$. Fixing $x_{<p}$ fixes the starting state of the walk (by fixing $x_{\pstart}$), the delay of the walk (by fixing $x_{\pdelay}$) and the states until time $t$ of the walk. 
	By the definition of $\bL(x) \in \bL_{j,t}$, we know that at the time $t$, the walk $W(x_{<p} \circ z)$ will also be at the state $j$. 
	Thus, for $\bL(x_{<p} \circ z)$ to be matched, we need the walk $W(x_{<p} \circ z)$ to reach the state $0$ or equivalently, we need the walk $W_z := w_t(=j),w_{t+1},\cdots,$ to reach $0$ starting from $j$. 
	 Note that the entire movement of the walk $W_z$ is determined by $z$.  
	
	On the other hand, consider the reflection walk $\barW(y_{<p} \circ z)$. Again, fixing $y_{<p}$ and the fact $\bR(y) \in \bR_{2m-j,t}$, in particular means that this walk is at state $2m-j$ at time $t$. 
	Consider the continuation of this walk as $\barW_z := \bar{w}_t (=2m-j), \bar{w}_{t+1},\cdots,$ starting from $2m-j$ which is fully determined by $z$. 
	Given $\barW_z$ is a reflection of $W_z$, if $W_z$ starting from $j$ reaches $0$, then, $\barW_z$ starting from $2m-j$ will reach $2m$. 
	But, that means that the walk $\barW(y_{<p} \circ z)$ ends in the state $2m$ which forces $\bR(y_{<p} \circ z)$ to never be matched in the blueprint. 
	
	Combining the above, we have that between the two vertices $\bL(x_{<p} \circ z)$ and $\bR(y_{<p} \circ z)$, at most one is matched in the blueprint, respecting the ban constraint. 
\end{proof}

	


We thus established that the blueprint $\bG$ is indeed a proper blueprint. 

\subsection{Value Calculation} 

The last step is to calculate the value of the blueprint which is $\card{\bE}/C^P$. 

\begin{lemma}\label{lem:value}
	The value of blueprint $\bG=\bG_m$ is 
	\[
		value(\bG) \geq \paren{1-\frac{1}{m}} \cdot \frac{2m}{3m+1}-3m \cdot 2^{-m}, 
	\]
\end{lemma}
\begin{proof}
	Given $\bE$ is a matching, and by construction, 
	\begin{align*}
		value(\bG) &= \frac{\card{\bE}}{C^P} \tag{by~\Cref{def:blueprint-value-approx}} \\
		&= \sum_{j=1}^{2m-1} \sum_{t = L}^{D} \frac{\min\paren{\card{\bL_{j,t}},\card{\bR_{2m-j,t}}}}{C^P} \tag{by the definition of edges} \\
		&\geq \sum_{j=1}^{2m-1} \sum_{t = L}^{D} \paren{\frac{\card{\bL_{j,t}}}{C^P} - 2^{-m} \cdot \frac{1}{D}} \tag{by~\Cref{lem:bjt-reflect}} \\
		&\geq \paren{\sum_{j=1}^{2m-1} \sum_{t = L}^{D} \frac{\card{\bL_{j,t}}}{C^P}} - 2m \cdot 2^{-m}. 
	\end{align*}
	We can now bound the inner sum separately as follows:
	\begin{align*}
		\sum_{j=1}^{2m-1} \sum_{t = L}^{D} \frac{\card{\bL_{j,t}}}{C^P} &= \sum_{j=1}^{2m-1} \sum_{t = L}^{D} \Pr\Paren{\text{$W$ $L$-restricted, $W$ reaches $0$, $\maxwalk{W}=j$, $\timemax{W} = t$}} \tag{by the definition of $\bL_{j,t}$} \\
		&= \sum_{j=1}^{2m-1} \Pr\Paren{\text{$W$ $L$-restricted, $W$ reaches $0$, $\maxwalk{W}=j$, $\timemax{W} \in [L:D]$}} \tag{as the events in the sum are all disjoint} \\
		&\geq \sum_{j=1}^{2m-1} \Pr\Paren{\text{$W$ $L$-restricted, $W$ reaches $0$, $\maxwalk{W}=j$, $d(x) \in [L:D-L]$}} \tag{since $W$ is $L$-restricted, having delay in $[L:D-L]$ ensures $\timemax{W} \in [L:D]$} \\
		&\geq \frac{D-2L}{D} \cdot \sum_{j=1}^{2m-1} \Pr\Paren{\text{$W$ $L$-restricted, $W$ reaches $0$, $\maxwalk{W}=j$}} \tag{as $d(x)$ is uniform independently on $[D]$ and $\card{[L:D-L]} = D-2L+1 > D-2L$} \\
		&= \frac{D-2L}{D} \cdot \Pr\Paren{\text{$W$ $L$-restricted, $W$ reaches $0$}} \tag{the sum is over all disjoint possibilities of reaching $0$} \\
		&\geq \frac{D-2L}{D} \cdot  \paren{\Pr\Paren{\text{$\widetilde W$ reaches $0$}}-2^{-m}} \tag{by~\Cref{clm:walk-length} and using $\widetilde W$ as a generic random walk} \\
		&= \frac{D-2L}{D} \cdot \paren{\frac{2m}{3m+1}-2^{-m}} \tag{by~\Cref{lem:start-state}} \\
		&\geq \paren{1-\frac{1}{m}} \cdot \frac{2m}{3m+1}-2^{-m} \tag{as $D = 2m \cdot L$}. 
	\end{align*}
	By plugging in this into the earlier bound, we have, 
	\[
		value(\bG) \geq \paren{1-\frac{1}{m}} \cdot \frac{2m}{3m+1}- 3m \cdot 2^{-m}, 
	\]
	concluding the proof. 
\end{proof}

The lower bound in \Cref{thm:blueprint} now follows immediately from~\Cref{lem:value} and the properness of the blueprint  established in~\Cref{clm:bG-matching,clm:bG-ban}. The upper bound is simply because 
\[
	value(\bG_m) \leq \Pr_{W \sim \dist_m}\paren{\text{$W$ reaches $0$}} = \frac{2m}{3m+1}, 
\]
by~\Cref{lem:start-state}, as every matched vertex in $\bL$ belongs to some $\bL_{j,t}$ and hence its walk has reached $0$. 

This concludes the proof of~\Cref{thm:blueprint}. 
